\chapter{Discussion and Analysis}
\label{ch:discussion}

\section{Interpretation of Results}

\subsection{Throughput Fairness}

The 3.2× throughput advantage for HIGH-priority devices confirms that weighted fair queuing effectively differentiates traffic. The weight ratio (HIGH=4, MEDIUM=2, LOW=1) translates to a 4:2:1 bandwidth distribution, closely matching the observed 5.8:3.2:1.8 Mbps allocation.

\textbf{Implication:} Software-based token bucket enforcement can achieve fairness comparable to hardware queuing systems.

\subsection{Latency and Jitter Reduction}

The 47\% latency reduction and 62\% jitter reduction for HIGH-priority traffic are substantial. These improvements are critical for real-time applications:
\begin{itemize}
    \item VoIP requires $<$ 150 ms latency and $<$ 30 ms jitter
    \item Video conferencing requires $<$ 200 ms latency and $<$ 50 ms jitter
    \item Online gaming requires $<$ 50 ms latency and $<$ 20 ms jitter
\end{itemize}

Our system brings latency from 68 ms (baseline) to 36 ms (QoS-enabled), meeting requirements for all three application classes.

\textbf{Implication:} Priority-based scheduling significantly improves user experience for latency-sensitive applications.

\subsection{Classification Accuracy}

The 92\% overall classification accuracy is acceptable for consumer-grade \qos. Misclassifications primarily occur for:
\begin{itemize}
    \item HTTPS traffic (port 443) used by multiple application types
    \item Applications using dynamic ports
    \item Peer-to-peer protocols
\end{itemize}

\textbf{Mitigation:} Future work should incorporate behavioral heuristics (packet size, inter-arrival time) to improve accuracy on encrypted traffic.

\subsection{Performance Overhead}

CPU usage (12\%) and memory footprint (68 MB peak) are well within acceptable bounds for mid-range Android devices. The sub-5ms per-packet latency overhead is negligible compared to typical cellular network latency (30–100 ms).

\textbf{Implication:} Userspace packet processing is viable for consumer-grade \qos on modern Android hardware.

\section{Comparison with Related Work}

\subsection{vs Kernel-Level Traffic Control}

Linux \texttt{tc} with HTB (Hierarchical Token Bucket) achieves $<$ 1 ms per-packet latency and sustains $>$ 100,000 pps. Our userspace implementation achieves 5 ms latency and 10,000 pps—sufficient for typical mobile hotspot scenarios (5–10 Mbps uplink, 500–1000 pps).

\textbf{Trade-off:} We sacrifice peak performance for portability (no root required).

\subsection{vs Commercial VPN Applications}

Commercial VPNs (NordVPN, ExpressVPN) focus on encryption and privacy, not \qos. Our system demonstrates that \vpn Service can be repurposed for bandwidth management.

\textbf{Novelty:} First consumer-grade \qos scheduler leveraging \vpn Service.

\subsection{vs SDN-Based Solutions}

SDN solutions (OpenFlow, P4) require infrastructure support unavailable in personal hotspot scenarios. Our system operates entirely on the end-user device.

\textbf{Advantage:} No infrastructure dependencies; works with any cellular carrier.

\section{Threats to Validity}

\subsection{Internal Validity}

\textbf{Threat:} iPerf3 synthetic traffic may not represent real application behavior.

\textbf{Mitigation:} Future work should include real application testing (Zoom, Netflix, online games).

\subsection{External Validity}

\textbf{Threat:} Laboratory settings may not reflect real-world usage patterns.

\textbf{Mitigation:} Field trials with diverse user populations are needed.

\subsection{Construct Validity}

\textbf{Threat:} \qos metrics (throughput, latency, jitter) may not fully capture user-perceived quality.

\textbf{Mitigation:} Future work should include subjective quality assessments (Mean Opinion Score).

\section{Limitations}

\subsection{Classification Limitations}

Port-based classification struggles with:
\begin{itemize}
    \item Encrypted traffic (HTTPS, VPN-over-VPN)
    \item Dynamic ports (peer-to-peer applications)
    \item Port obfuscation (intentional misuse of well-known ports)
\end{itemize}

\subsection{Uplink Bandwidth Estimation}

Manual uplink configuration is error-prone. Automatic estimation is challenging due to variable cellular link capacity.

\subsection{MAC Address Resolution}

Current implementation does not resolve MAC addresses, limiting priority persistence across DHCP lease renewals.

\section{Lessons Learned}

\subsection{Technical Lessons}

\begin{enumerate}
    \item \textbf{Userspace performance is sufficient:} Contrary to conventional wisdom, userspace packet processing achieves acceptable performance for consumer-grade \qos.
    
    \item \textbf{Simple heuristics work well:} Port-based classification achieves 92\% accuracy despite its simplicity.
    
    \item \textbf{Reactive architecture simplifies state management:} Kotlin StateFlow eliminates entire classes of synchronization bugs.
\end{enumerate}

\subsection{Research Methodology Lessons}

\begin{enumerate}
    \item \textbf{Design science is effective:} Iterating between design, implementation, and evaluation enables rapid refinement.
    
    \item \textbf{Open-source accelerates validation:} Releasing the implementation enables independent replication.
    
    \item \textbf{Controlled experiments are necessary but insufficient:} Field trials are essential for real-world validation.
\end{enumerate}

\section{Implications for Practice}

\subsection{For End Users}

\begin{itemize}
    \item Prioritize video calls over background downloads
    \item Ensure online gaming remains responsive during concurrent usage
    \item Manage bandwidth in multi-user households
\end{itemize}

\subsection{For Developers}

\begin{itemize}
    \item \vpn Service can be repurposed for network control beyond encryption
    \item Userspace packet processing is viable on modern Android hardware
    \item Reactive architecture patterns simplify state management
\end{itemize}

\subsection{For Researchers}

\begin{itemize}
    \item Consumer-grade \qos is feasible without root access
    \item Port-based classification remains effective for common traffic
    \item Software token buckets can achieve acceptable performance
\end{itemize}
